
\section{Appendix A. Mathematical Derivations}

\subsection{A.1 Derivation of $\alpha$-Entmax Closed Form}

\textbf{Problem.} Solve:

\[
p^* = \arg\max_{p \in \Delta^{d-1}} \langle p, s \rangle + H_\alpha^T(p)
\tag{A.1}
\]

where $H_\alpha^T(p) = \frac{1}{\alpha(\alpha-1)} \sum_i (p_i - p_i^\alpha)$.

\textbf{Lagrangian.}

\[
\mathcal{L}(p, \tau) = \langle p, s \rangle + \frac{1}{\alpha(\alpha-1)} \sum_i (p_i - p_i^\alpha) - \tau \left(\sum_i p_i - 1\right)
\tag{A.2}
\]

\textbf{Stationarity.} $\partial \mathcal{L} / \partial p_i = 0$:

\[
s_i + \frac{1 - \alpha p_i^{\alpha-1}}{\alpha(\alpha-1)} - \tau = 0
\tag{A.3}
\]

Rearranging:

\[
p_i^{\alpha - 1} = (\alpha - 1) s_i - \tau + \frac{1}{\alpha}
\tag{A.4}
\]

Let $\tilde\tau = \tau - 1/\alpha$. Then:

\[
p_i = \max\left(0, (\alpha - 1) s_i - \tilde\tau\right)^{1/(\alpha - 1)}
\tag{A.5}
\]

with $\tilde\tau$ chosen so that $\sum_i p_i = 1$. $\square$

\subsection{A.2 Bisection Convergence}

\textbf{Claim.} Bisection on $[\tau_\text{lo}, \tau_\text{hi}]$ converges to 
precision $\epsilon$ in $O(\log((\tau_\text{hi} - \tau_\text{lo})/\epsilon))$ 
iterations.

\textbf{Proof.} Standard bisection result. With initial range $O(1)$ and 25 
iterations, precision is $O(2^{-25}) = O(10^{-7.5})$. $\square$

\subsection{A.3 Support Sparsity Bound}

\textbf{Claim.} At convergence, the support of $\alpha$-entmax is:
\[
S = \{i : (\alpha - 1) s_i > \tau^*\}
\]

\textbf{Proof.} Direct from the closed form: $p_i > 0$ iff $(\alpha - 1) s_i > \tau^*$. $\square$

\subsection{A.4 Attention Sink Preservation}

\textbf{Claim.} For causal decoders with splash masking, preserving first $S$ 
tokens does not violate causality.

\textbf{Proof.} For query $i$, attention mask includes key $j$ if:
\begin{itemize}
\item $j < S$ (sink) OR
\item $0 \leq i - j \leq W$ (local window)

\end{itemize}
Case 1: Sink. If $j < S \leq i$ (assuming $i \geq S$), then $j \leq i$, 
so causality is preserved. For $i < S$, only $j \leq i$ are included 
because $i - j \geq 0$ requires $j \leq i$.

Case 2: Local. $0 \leq i - j \leq W$ implies $j \leq i$. Causality holds.

Therefore, all attended keys satisfy $j \leq i$. $\square$

\subsection{A.5 Parameter Count Analysis}

For square $d \times d$ projection:

\[
\text{Original params} = d^2
\tag{A.6}
\]

\[
\text{AdaSplash params} = d^2 + r_L (d + d) = d^2 + 2 r_L d
\tag{A.7}
\]

\textbf{Change:} $+2 r_L d$ per projection.

Across 4 projections: $+8 r_L d$ per layer.

Across $N_l$ layers: $+8 r_L d N_l$.

For DistilBERT ($d = 768$, $r_L = 16$, $N_l = 6$): $+589{,}824$ params.

\subsection{A.6 Entmax vs Softmax Sparsity}

\textbf{Claim.} For scores with high variance, $\alpha$-entmax produces more zeros 
than softmax.

\textbf{Proof sketch.} Softmax assigns $p_i > 0$ for all $i$. Entmax assigns 
$p_i = 0$ when $(\alpha - 1) s_i < \tau^\textit{$. The threshold $\tau^}$ scales 
with the score distribution; for high variance, fewer scores exceed 
$\tau^*$, giving more sparsity. $\square$

\subsection{A.7 Softmax Fallback}

\textbf{Claim.} When support is empty, the softmax fallback is a valid probability 
distribution.

\textbf{Proof.} Softmax always produces valid probabilities. $\square$

\subsection{A.8 Causality of Entmax}

\textbf{Claim.} $\alpha$-entmax preserves causality under causal masking.

\textbf{Proof.} Entmax is a row-wise operation on the score matrix. Masking with 
$-\infty$ ensures that masked entries receive zero weight in the output. 
Row-wise independence preserves row-specific causality. $\square$